\documentclass[10pt,a4paper]{article}
\usepackage[T1]{fontenc}
\usepackage{lmodern}
\usepackage{amsmath,amssymb,mathtools}
\usepackage{geometry}
\usepackage{microtype}
\usepackage{enumitem}
\usepackage{booktabs}
\usepackage{xcolor}
\usepackage{hyperref}
\usepackage{titlesec}
\usepackage{fancyhdr}
\geometry{top=20mm,bottom=19mm,left=21mm,right=21mm}
\definecolor{navy}{HTML}{17324D}
\definecolor{accent}{HTML}{126F78}
\definecolor{muted}{HTML}{526579}
\hypersetup{colorlinks=true,urlcolor=accent,linkcolor=accent}
\titleformat{\section}{\color{navy}\normalfont\bfseries\large}{}{0pt}{}
\titleformat{\subsection}{\color{navy}\normalfont\bfseries\normalsize}{}{0pt}{}
\titlespacing*{\section}{0pt}{9pt}{3pt}
\titlespacing*{\subsection}{0pt}{6pt}{2pt}
\setlength{\parindent}{0pt}
\setlength{\parskip}{4pt}
\setlist[enumerate]{leftmargin=*,topsep=2pt,itemsep=2pt,parsep=0pt}
\setlist[itemize]{leftmargin=*,topsep=2pt,itemsep=1.5pt,parsep=0pt}
\allowdisplaybreaks[1]
\pagestyle{fancy}
\fancyhf{}
\fancyhead[L]{\footnotesize\color{muted}Research note | Network-size detection}
\fancyhead[R]{\footnotesize\color{muted}Haehne et al. (2019)}
\fancyfoot[C]{\footnotesize\thepage\ / 2}
\renewcommand{\headrulewidth}{0.35pt}
\setlength{\headheight}{13pt}
\begin{document}
\begin{center}
{\Large\bfseries\color{navy} Detecting Hidden Units and Network Size From Perceptible Dynamics}\\[2pt]
{\small H. Haehne, J. Casadiego, J. Peinke, and M. Timme,\quad
\emph{Physical Review Letters} \textbf{122}, 158301 (2019)}\\[-1pt]
{\small \href{https://doi.org/10.1103/PhysRevLett.122.158301}{doi:10.1103/PhysRevLett.122.158301}}
\end{center}
\vspace{-3pt}
\noindent\colorbox{navy!6}{\parbox{\dimexpr\linewidth-2\fboxsep\relax}{\textbf{Main question.}
A network has $N$ interacting dynamical units, but only $n<N$ are measured. Can we infer the total size $N$---and hence the $h:=N-n$ hidden units---without knowing the governing equations or network connections?\\[2pt]
\textbf{Main idea.} Arrange observations from \emph{several distinct transient trajectories} into a detection matrix $T$. Under sufficient excitation and observability, its rank equals the full system size: $\boxed{\operatorname{rank}(T)=N}$.
\\
\textbf{Sufficient Excitation}: it means if we don not have enough different experiments or experiments have the same response then we can not detect the hidden units.
\\
\textbf{Observability}: it means if the hidden units do not influence the measured units then we can not detect the hidden units.
}}

\section{1. Notation and dynamical model}
\begin{center}
\small
\begin{tabular}{@{}ll@{} }
\toprule
$N$ / $n$ & Total number of units / number of directly measured units ($n<N$) \\
$M$ / $k$ & Number of distinct trajectories (experiments) / sampling times per trajectory \\
$\mathbf z(t)\in\mathbb R^N$ & Full network state (including unmeasured variables) \\
$\mathbf y^{(m)}(t)\in\mathbb R^n$ & Measurements from experiment $m=1,\ldots,M$ \\
$T\in\mathbb R^{kn\times M}$ & Detection matrix assembled from measured trajectories \\
\bottomrule
\end{tabular}
\end{center}
The unknown network evolves according to
\begin{equation}
  \dot{\mathbf z}(t)=\mathbf F(\mathbf z(t)). \tag{1}
\end{equation}
Near a fixed point $\mathbf z^*$, where $\mathbf F(\mathbf z^*)=0$, let
$\mathbf x=\mathbf z-\mathbf z^*$. A first-order approximation gives
\begin{equation}
  \dot{\mathbf x}=A\mathbf x,\quad
  A= D\mathbf F(\mathbf z^*),\quad
  \mathbf x(t)=e^{At}\mathbf x(0). \tag{2}
\end{equation}
Here $A$ is the (unknown) Jacobian, which describes local influences between units.
Only $n$ states are observed. With the selection matrix
$C=[\,I_n\;0\,]\in\mathbb R^{n\times N}$,
\begin{equation}
  \mathbf y^{(m)}(t)=C e^{At}\mathbf x^{(m)}(0)
  =\Theta(t)\mathbf x^{(m)}(0),\qquad \Theta(t)=Ce^{At}. \tag{3}
\end{equation}
The key point is that an observed response can depend on \emph{all} $N$ initial-state components, including those that are hidden.

\section{2. Construct the detection matrix and derive the rank bound}
At each common time $t_j$ measured relative to the start of the experiment, arrange the $M$ responses as columns:
\[
 Y(t_j)=[\mathbf y^{(1)}(t_j)\;\cdots\;\mathbf y^{(M)}(t_j)]
 \in\mathbb R^{n\times M}.
\]
Stack the measurements at $k$ times, preserving the same experiment in each column:
\begin{equation}
\boxed{T=\begin{bmatrix}Y(t_1)\\\vdots\\Y(t_k)\end{bmatrix}}
=\underbrace{\begin{bmatrix}\Theta(t_1)\\\vdots\\\Theta(t_k)\end{bmatrix}}_{\mathcal O_k\in\mathbb R^{kn\times N}}
\underbrace{[\mathbf x^{(1)}(0)\;\cdots\;\mathbf x^{(M)}(0)]}_{X_0\in\mathbb R^{N\times M}}.
\tag{4}
\end{equation}
By the rank inequality for a matrix product,
\begin{equation}
\operatorname{rank}(T)\le\min\{\operatorname{rank}(\mathcal O_k),
\operatorname{rank}(X_0)\}\le\min\{kn,N,M\}\le N.\tag{5}
\end{equation}
If $\mathcal O_k$ has \emph{full column rank} $N$ (the chosen measurements reveal all state directions) and $X_0$ has \emph{full row rank} $N$ (the initial conditions span them), then
\begin{equation}
\boxed{\operatorname{rank}(T)=N,\qquad \widehat h=N-n=\operatorname{rank}(T)-n.}\tag{6}
\end{equation}
For generic informative data, the paper uses $kn>N$ and $M>N$; these size inequalities alone do \emph{not} guarantee full rank.

\newpage
\section{3. Practical algorithm: estimate rank using singular values}
In measured data, noise makes exact numerical rank unreliable. The paper therefore computes the ordered singular values of $T$,
\[
 T=U\Sigma V^\top,\qquad
 \sigma_1\ge\sigma_2\ge\cdots\ge\sigma_b\ge0,\quad b=\min(kn,M),
\]
and looks for an unusually large gap between successive values on a logarithmic scale:
\begin{equation}
\boxed{\Delta_j=\log\sigma_j-\log\sigma_{j+1}
=\log\left(\frac{\sigma_j}{\sigma_{j+1}}\right),\quad j=1,\ldots,b-1.}\tag{7}
\end{equation}
The estimated network size is the \emph{index} of the relevant gap, not its numerical height. The full iterative procedure is:
\begin{enumerate}[label=\textbf{\arabic*.}]
\item \textbf{Collect trials:} Start with $M=n+1$ distinct transient trajectories, measured from the same $n$ accessible units.
\item \textbf{Sample each trial:} Choose $k=\lceil M/n\rceil$ equally separated relative sampling times on each trajectory; construct $T$ as in (4).
\item \textbf{Estimate a gap index:} Compute the singular values and $\Delta_j$; among $j\ge n$, omit the potentially misleading indices $\{n,2n,\ldots\}$ and $\{n+1,2n+1,\ldots\}$. Identify the index of the largest remaining gap.
\item \textbf{Check robustness:} Repeat the sampling and gap calculation after retaining only $n-1,n-2,n-3,n-4$ measured units (when these counts are positive). Take the \emph{median of the estimated indices} as $\widehat N^{(M)}$.
\item \textbf{Increase experiments:} Add another trajectory ($M\leftarrow M+1$) and repeat until $\widehat N^{(M)}$ stops increasing. Report $\widehat N=\widehat N^{(M)}$ and $\widehat h=\widehat N-n$.
\end{enumerate}
\vspace{-2pt}
\textbf{Intuition:} Increasing $M$ and $k$ can reveal new independent directions. Once all $N$ directions are represented, the rank estimate saturates; a spectral gap appears after $\sigma_N$. A noisy or temporarily flat estimate, however, is not by itself proof of saturation.

\section{4. Assumptions and scope}
\begin{itemize}
\item \textbf{Enough independent trials:} In the paper's generic full-rank argument, $M>N$ and $kn>N$. Distinct trials must excite enough independent initial-state directions; repeated identical trials are not sufficient.
\item \textbf{Observability through measured units:} Hidden states must leave distinguishable signatures in the observed responses. A dynamically invisible unit cannot be counted by this procedure.
\item \textbf{Comparable transient measurements:} Sampling times are aligned relative to each trial's start, and the underlying dynamics are consistent across trials.
\item \textbf{Valid local linearization:} Near fixed points, (2) is a suitable approximation. For nonlinear periodic or chaotic motion, the method instead uses differences of \emph{nearby trajectories}, $\delta\mathbf z(t)\simeq D\Phi_{t-t^*}|_{\mathbf z^*}\,\delta\mathbf z(t^*)$, so that a comparable linear relation holds locally. Trajectories must remain sufficiently close for this approximation.
\item \textbf{Noise and conditioning:} Measurement noise, weak coupling, and poorly separated singular values can cause over- or underestimation. The exact-rank theorem does not guarantee a correct rank estimate on every finite, noisy dataset.
\end{itemize}

\section{5. What did the authors find?}
The paper demonstrates recovery of total network size in numerical experiments with (i) linear diffusively coupled systems near fixed points, (ii) nonlinear phase-locked Kuramoto oscillators, (iii) periodic and (iv) chaotic R\"ossler networks, and (v) nonlinear heterogeneous Michaelis--Menten biochemical networks with noise. In favorable examples, size detection worked when only about \emph{10\% of units were measured}. For tested linear networks, the minimum number of trials was consistent with $M_{\min}=N+1$; in the shown Kuramoto experiments it was $N+2$ because one trial served as a reference. More trials can allow accurate detection in noisy settings.

\noindent\colorbox{navy!6}{\parbox{\dimexpr\linewidth-2\fboxsep\relax}{\textbf{Final result and limitation.} The method can infer the \emph{number} of hidden units from sufficiently informative transient observations without first identifying the full governing equations or connectivity. It does \emph{not} reconstruct hidden states, reveal where the hidden units are, or identify their exact connections.}}

\vspace{3pt}
\color{muted}\textbf{Many Parts of this summary are generated by ChatGPT.}
\\
{\footnotesize\color{muted}\textbf{Source:} Haehne \emph{et al.}, \emph{Phys. Rev. Lett.} \textbf{122}, 158301 (2019), especially Eqs.~(1)--(13), Figs.~1--4 and the concluding discussion.}
\end{document}
